\documentclass[letterpaper,12pt]{article}
%\usepackage{harvard}
%\usepackage{subfigure}
%\usepackage[active]{srcltx}
\usepackage{graphicx}
\usepackage{pdflscape}
\usepackage{amsmath}
\usepackage{amsthm}
\usepackage{lscape}
\usepackage{enumerate}
\usepackage{float}
\usepackage{grffile}
\usepackage{enumitem}
\usepackage{relsize}
\usepackage{bm}
\usepackage{tablefootnote}
\usepackage{footnote}
\usepackage{sectsty}
\usepackage{xcolor}
\usepackage{caption}
\usepackage{color}
\usepackage{natbib}
\usepackage{multirow}
\usepackage{bbm}
\usepackage[T1]{fontenc}
\usepackage[multiple]{footmisc}
\usepackage{appendix}
\usepackage{times}
\usepackage{relsize}
\usepackage{caption}
\usepackage{newtxtext,newtxmath}
\usepackage{enumitem}


%%%%%%%%%%%%%%%%%%%%%%%%%%%%%%%%%%%%%%%%%%%%%%%%%%%%%%%%%%%%
%%%   Title Page Information
%%%%%%%%%%%%%%%%%%%%%%%%%%%%%%%%%%%%%%%%%%%%%%%%%%%%%%%%%%%%
\title{\vspace{-1.2in} ECON712: Midterm Exam}
\author{Leandro Freylejer}
\date{February 26th 2026}

%%%%%%%%%%%%%%%%%%%%%%%%%%%%%%%%%%%%%%%%%%%%%%%%%%%%%%%%%%%%
%%%   Double and Single Space Commands
%%%%%%%%%%%%%%%%%%%%%%%%%%%%%%%%%%%%%%%%%%%%%%%%%%%%%%%%%%%%
\newlength{\defbaselineskip}
\setlength{\defbaselineskip}{\baselineskip}
\newcommand{\setlinespacing}[1]{\setlength{\baselineskip}{#1 \defbaselineskip}}
\newcommand{\doublespacing}{\setlength{\baselineskip}{1.3 \defbaselineskip}}
\newcommand{\singlespacing}{\setlength{\baselineskip}{1\defbaselineskip}}
\renewcommand{\thesubsection}{\arabic{subsection}}
\newtheoremstyle{dotlessS}{}{}{\itshape}{}{\color{dg}\bfseries}{}{ }{}
\theoremstyle{dotlessS}
\newtheorem{theorem}{Theorem}
\newtheorem{lemma}{Lemma}
\newtheorem{proposition}{Proposition}
\newtheorem{assumption}[theorem]{Assumption}
\newtheorem{corollary}{Corollary}
\newtheorem{definition}{Definition}
\newenvironment{example}[1][Example]{\begin{trivlist}
\item[\hskip \labelsep {\bfseries #1}]}{\end{trivlist}}
\newenvironment{remark}[1][Remark]{\begin{trivlist}
\item[\hskip \labelsep {\bfseries #1}]}{\end{trivlist}}

\newcommand\Pitem{%
  \addtocounter{enumi}{-1}%
  \renewcommand\theenumi{\arabic{\enumi}'}%
  \item%
  \renewcommand\theenumi{\arabic{enumi}}%
}

%%%%%%%%%%%%%%%%%%%%%%%%%%%%%%%%%%%%%%%%%%%%%%%%%%%%%%%%%%%%
%%%   Set Margins Here
%%%%%%%%%%%%%%%%%%%%%%%%%%%%%%%%%%%%%%%%%%%%%%%%%%%%%%%%%
\textwidth=6.5in
\textheight=9in
\oddsidemargin=0.10in
\evensidemargin=0.10in
\topmargin=-0.5in
%\parskip=3pt

%%%%%%%%%%%%%%%%%%%%%%%%%%%%%%%%%%%%%%%%%%%%%%%%%%%%%%%%%%%%%
%%%   Actual Document Begins Here
%%%%%%%%%%%%%%%%%%%%%%%%%%%%%%%%%%%%%%%%%%%%%%%%%%%%%%%%%%%%%
\begin{document}
\pagenumbering{arabic}
\maketitle
\vspace{-.2in}
\noindent There are five questions in this exam for a total of 53 points. Good Luck!!!

\section*{Question I: Linear transformation of variables (8 points)}
Assume the regression model:
\begin{eqnarray*}
\bm{y}=\bm{\iota_n} \: \hat{\gamma} + \bm{X} \: \bm{\hat{\beta}}+\bm{\hat{u}}
\end{eqnarray*}
where
\begin{itemize}
\item $\bm{y}$ is a $nx1$ vector of dependent variables
\item $\bm{\iota_n}$ is an $nx1$ vector of ones, 
\item $\bm{X}$ is a $nxk$ matrix of regressors, 
\item $\hat{\gamma}$ is a scalar OLS/MM estimate
\item $\bm{\hat{\beta}}$ is a $kx1$ vector of OLS/MM estimates
\item $\bm{\hat{u}}$ is a $nx1$ vector of residuals
\end{itemize}
Define the following transformed variables:
\begin{eqnarray*}
\bm{z}=\bm{y}- \bm{\iota_n} \overline{y}
\end{eqnarray*}
and
\begin{eqnarray*}
\bm{W}=\bm{X}-\bm{\iota_n} \bm{\overline{x}}^T
\end{eqnarray*}
with
\begin{eqnarray*}
\begin{array}{lllll}
\overline{y}=\frac{1}{n}\underset{i=1}{\overset{n}{\sum}}y_{i} &,& \bm{\overline{x}}=\begin{pmatrix} \bar{x}_1 \\ \bar{x}_2 \\ \vdots  \\  \bar{x}_k \end{pmatrix} &,& \bar{x}_j=\frac{1}{n}\underset{i=1}{\overset{n}{\sum}}x_{ij} 
\end{array}
\end{eqnarray*}
Suppose that a researcher uses OLS/MM to estimate the following regression isntead:
\begin{eqnarray*}
\bm{z}= \bm{W} \: \bm{\hat{\alpha}}+\bm{\hat{\epsilon}}
\end{eqnarray*}
Derive the relationship between $\bm{\hat{\alpha}}$ and $\bm{\hat{\beta}}$. Are the residuals from these two regressions equal? 


\section*{Question II: Finite Sample Properties of OLS (8 points)}
Assume that the model $\bm{y}=\bm{X} \: \bm{\beta}+\bm{u}$ satisfies the Gauss-Markov assumptions (i.e. classical assumptions), where $\bm{X}$ is a $n x k$ matrix and $\bm{y}$ is a $n x 1$ vector of observations. Define a new matrix $\bm{Z}=a \: \bm{X}$, where $a>0$ is a scalar and a new vector $\bm{s}=b \: \bm{y}$, where $b>0$ is a scalar. Suppose a researcher draws a random sample $n$ and run the following regression:
\begin{eqnarray*}
\bm{s}=\bm{Z} \bm{\hat{\beta}}+\bm{\hat{u}}
\end{eqnarray*}
Prove whether $\bm{\hat{\beta}}$ an unbiased estimator of $\bm{\beta}$? [Hint: you need to consider two separate cases: $a \neq b$ and $a=b$] 



\section*{Question III: Linear regression and selection bias }
Suppose potential outcomes are defined as:
\begin{eqnarray*}
y_i=y_{0i}+(y_{1i}-y_{0i}) \: D_i,
\end{eqnarray*}
where $D_i \in \left\{0,1\right\}$ is a treatment dummy. Assume that the econometrician runs the following \textit{population} regression (i.e. solves for $\beta_0$ and $\beta_1$ using OLS/MM):
\begin{eqnarray*}
y_i=\beta_0+\beta_1 \: D_i+u_i
\end{eqnarray*}
Using this information answer the following questions:
\begin{enumerate}[label=(\alph*)]
\item (5 points) Solve for $\beta_1$ in terms of conditional expectations (show your work)
\item (5 points) Suppose that there are homogeneous treatment effects: $y_{1i}-y_{0i}=\rho$ , show that $\beta_1$ is equal to the causal effect of the treatment, defined as the potential outcome, plus a term denoting selection bias. 
\item (5 points) Suppose now that there are heterogeneous treatment effects: $y_{1i}-y_{0i}=\rho_i$ , show that $\beta_1$ is equal to the average treatment effect on all individuals, plus a term that captures the additional effect on treated individuals, plus selection bias [Hint: to answer this question you \textit{only} need to add and subtract $\mathbb{E}[\rho_i]$, and then reorganize terms]
 
\end{enumerate}
     
\section*{Question IV: Goodness of Fit (5 points)}
Using projection matrices the Pythagorean theorem can be expressed as:
\begin{eqnarray*}
\|\bm{y}\|^2=\|\bm{P_X \: y}\|^2+\|\bm{M_X \: y}\|^2
\end{eqnarray*}
This is the basis from which the following two versions of the $R^2$ are obtained:
\begin{eqnarray*}
\begin{array}{lll}
R_u^2=\frac{\|\bm{P_X \: y}\|^2}{\|\bm{y}\|^2} &, & R_c^2=\frac{\|\bm{P_X \: M_{\iota} \: y}\|^2}{\|\bm{M_{\iota} \:y}\|^2}
\end{array}
\end{eqnarray*}
Assume that $\bm{X}$ includes the constant term. Consider the transformed variable $\bm{z}=\bm{y}-c \: \bm{\iota}$, prove that this transformation changes the $R_u^2$, but leaves  $R_c^2$  


\section*{Question V: Conditional Expectation Function }

Suppose that the random variables $Y$ and $X$ only take the value $0$ and $1$ and they have the following joint probability distribution: 
\begin{table}[H]
\centering
\resizebox{.25\textwidth}{!}{%
\begin{tabular}{|c| c |c|}
\hline
         & Y=0 & Y=1 \\
         \hline
X=0  & 0.2 & 0.6 \\
X=1  & 0.1 & 0.1 \\
\hline
\end{tabular}}
\label{tab:three_by_three}
\end{table}

Using this information, answer the following questions
\begin{enumerate}[label=(\alph*)]
\item (7 points) Compute $\mathbb{E}[Y|X=0]$ and $\mathbb{E}[Y|X=1]$
\item (5 points) Derive the linear conditional expectation function: $m(X)=\mathbb{E}[Y|X]$  
\item (5 points) Define $u=Y-m(X)$ and show $\mathbb{E}[u|X]=0$. Use the LIE to show that $Cov(X,u)=0$ [Hint: show that $u$ is mean independent]
\end{enumerate}


\end{document}
